\documentclass[11pt]{article}
\usepackage[a4paper,margin=27mm]{geometry}
\usepackage{amsmath,amssymb,amsthm,hyperref}
\newcommand{\Jac}{\operatorname{Jac}}
\newcommand{\ord}{\operatorname{ord}}
\newcommand{\J}{\mathcal J}
\newtheorem{theorem}{Theorem}
\newtheorem{lemma}{Lemma}
\title{A local differential proof of the Rees discrepancy estimate}
\author{}
\date{6 October 2026}
\begin{document}
\maketitle

\begin{abstract}
For a coherent ideal on a complex manifold, a finite normalization map
and the differentials of its generator ratios give a discrepancy bound
at a Rees divisor. Combining it with the local Jacobian--multiplier
inclusion recovers the estimate needed in special K\"ahler Fujita
approximation. The proof does not assume that the blowup differential
is injective in the normal direction. This is an unreviewed candidate,
with no original problem or new scoped paper counted as solved.
\end{abstract}

\section{Precise local statement}

Let $X$ be a complex manifold of dimension $n\geq1$, and let
$\mathfrak a\ne0$ be a coherent ideal. Let
$p:Y\to X$ be its normalized blowup. Normality makes $Y$ smooth
at the generic point of every prime divisor. For such a divisor
$E$ write $v=\ord_E$. The coefficient
\[
 k_E=\ord_E K_{Y/X}
\]
means the order of the pullback of a local nonvanishing top form
on $X$, measured against a local nonvanishing top form at the
generic point of $E$. No $\mathbf Q$-Cartier assumption on
$K_Y$ is needed to define this coefficient.

For local holomorphic generators $f_0,\ldots,f_N$ of an ideal
$\mathfrak b$, set
\[
 \Jac_1(\mathfrak b)=
 \left(f_j,\frac{\partial f_j}{\partial x_i}:
              0\leq j\leq N,\ 1\leq i\leq n\right).
 \tag{1}
\]
Changing holomorphic coordinates or generators leaves this
ideal unchanged: the chain rule multiplies derivatives by an
invertible holomorphic matrix, and derivatives of coefficient
functions add multiples of the generators.
Use the ordinary exponent-one analytic multiplier ideal
$\J(\mathfrak b)$.

\begin{theorem}
For every prime divisor $E$ on $Y$,
\[
 k_E\leq n\{v(\mathfrak a)-v(\J(\mathfrak a))\}.
 \tag{2}
\]
For a $p$-exceptional divisor the following stronger bound holds:
\[
 k_E+1\leq n\{v(\mathfrak a)-v(\J(\mathfrak a))\}.
 \tag{3}
\]
The assertion is local on $X$ and requires neither projectivity
nor a global K\"ahler form.
\end{theorem}

The key new calculation is the following lemma.

\begin{lemma}
Suppose that $E$ is a prime divisor of the normalized blowup of
$\mathfrak b$ with $m=v(\mathfrak b)>0$. Put
$\ell=v(\Jac_1(\mathfrak b))$. Then
\[
 k_E+1\leq n(m-\ell).
 \tag{4}
\]
\end{lemma}

\begin{proof}
Choose a generator $f_0$ of order $m$ along $E$ which generates
$\mathfrak b\mathcal O_Y$ near its generic point. This is possible
on one of the usual blowup charts. The functions
$r_j=p^*(f_j/f_0)$ are holomorphic there.
The normalized blowup maps finitely to the ordinary blowup,
which on this chart embeds in the product of a coordinate
neighborhood of $X$ with an affine space of generator ratios.

Consequently, at a general smooth point of $E$, the restrictions
of the base coordinate functions $p^*x_i$ and of the ratios
$r_j$ have differential rank $n-1$ on $E$.
Here is the rank justification: the finite map restricted to
$E$ has an $(n-1)$-dimensional image; in characteristic zero
it has full rank at a general smooth point of that image.
The embedding by coordinate functions and ratios then supplies
$n-1$ independent restricted differentials. Pick $u_1,\ldots,
u_{n-1}$ from these coordinate functions and ratios with
independent restrictions.

Take a local equation $s=0$ for $E$ and write
$p^*f_0=s^m h$ with $h$ a holomorphic unit. Since the restricted
$du_i$ are independent, the top form
\[
 d(p^*f_0)\wedge du_1\wedge\cdots\wedge du_{n-1}
\]
has exact order $m-1$ along $E$. Indeed,
$d(s^mh)=m s^{m-1}h\,ds+s^m dh$, and the first term has a
nonzero tangential wedge on $E$.

Compute the same form on the base outside the ideal locus.
Every coefficient of $df_0$ in the $dx_i$ basis has valuation
at least $\ell$. A ratio differential satisfies
\[
 d(f_j/f_0)=\frac{f_0\,df_j-f_j\,df_0}{f_0^2};
 \tag{5}
\]
its coefficients have valuation at least $\ell-m$.
A coordinate differential has coefficients of valuation
zero or infinite, hence also at least $\ell-m$, because
$\Jac_1(\mathfrak b)$ contains $\mathfrak b$ and
therefore $\ell\leq m$.

The determinant expressing the base top form in
$dx_1\wedge\cdots\wedge dx_n$ thus has valuation at least
\[
 \ell+(n-1)(\ell-m)=m-n(m-\ell).
 \tag{6}
\]
Pullback of the base volume form contributes order $k_E$.
Comparison with the exact order $m-1$ gives
\[
 m-1\geq k_E+m-n(m-\ell),
\]
which is (4).
All valuation statements concern the generic divisor and are
unchanged after shrinking the coordinate neighborhoods.
\end{proof}

\section{The multiplier-ideal step and divisorial factors}

We use the exponent-one Jacobian--multiplier relation from
Ein--Lazarsfeld--Smith--Varolin, Theorem4.2(i) with $m=1$.
Its required codimension condition holds for ideals with no
divisorial factor. For clarity its analytic-local specialization
can be checked as follows.

\begin{lemma}
If $\operatorname{codim}V(\mathfrak b)\geq2$, then
$\Jac_1(\mathfrak b)\subset\J(\mathfrak b)$.
\end{lemma}

\begin{proof}
On a log resolution let $P$ have ideal order $q$ and canonical
order $b$. Only $q-b>0$ needs checking. Its center has codimension
$c\geq2$. Choose coordinates $x_1,\ldots,x_c$ vanishing on
that center. Near a generic point of $P=\{s=0\}$, generators
pull back as $s^q h$, and these $x_i$ as $s h_i$.
In the pullback of
$df_j\wedge\bigwedge_{i\ne l}dx_i$, at most one factor can use
$ds$. Its order is therefore at least $q+c-2\geq q$.
The same order is
$v_P(\partial f_j/\partial x_l)+b$. Thus every derivative
satisfies the multiplier criterion
$v_P(\partial f_j/\partial x_l)\geq q-b$.
Generators do too, since $\mathfrak b\subset\J(\mathfrak b)$.
The defining log-resolution formula proves the inclusion.
\end{proof}

\begin{proof}[Proof of the theorem]
Work locally and factor
$\mathfrak a=\mathfrak b\mathcal O_X(-G)$, where $G$ is
effective integral and $\mathfrak b$ has no divisorial factor.
Multiplying an ideal by an invertible ideal leaves its blowup,
and hence its normalized blowup, unchanged. The defining
log-resolution formula also gives
\[
 \J(\mathfrak a)=\J(\mathfrak b)\mathcal O_X(-G).
 \tag{7}
\]
It follows that
$v(\mathfrak a)-v(\J(\mathfrak a))
 =v(\mathfrak b)-v(\J(\mathfrak b))$.

For a $p$-exceptional divisor $E$, its center is in the
nontrivial ideal locus of $\mathfrak b$; hence
$m=v(\mathfrak b)>0$. Set $\ell=v(\Jac_1(\mathfrak b))$.
The last lemma gives $\ell\geq v(\J(\mathfrak b))$.
Equation (4) then yields
\[
 k_E+1\leq n(m-\ell)
 \leq n\{v(\mathfrak b)-v(\J(\mathfrak b))\},
\]
which is (3) and implies (2).
For a nonexceptional prime divisor, the birational map is an
isomorphism generically over a prime divisor of the smooth
base, so $k_E=0$. Also
$\mathfrak a\subset\J(\mathfrak a)$, giving a nonnegative
right-hand side in (2). This treats every prime divisor.
In dimension one the blowup is locally the identity and this
last case is the whole assertion.
\end{proof}

\section{Application and verification boundary}

The analytic discrepancy input in Trusiani2609.31089v1,
Proposition4.1, concerns precisely the local estimate (2).
Subject to independent verification of this candidate, the
proof above supplies that input on arbitrary smooth complex
manifolds. In particular it avoids any assertion that the
differential of a normalized blowup is nonzero on its normal
line at a generic exceptional point. It neither assumes a
global ample line bundle nor transfers algebraic Nadel
vanishing into the nonprojective setting.

This does not by itself prove the full special Fujita theorem.
The regularization construction, positivity estimates and
orthogonality limits in that theorem remain separate inputs
to the entropy-recovery programme. Nor does it prove ray
entropy recovery, a maximal-ray slope equality, equivariant
extensions, or degenerating-family compactness. Original P19.1
and P19.2--P19.4 remain open.

The background used here is local resolution of coherent
ideals, the log-resolution definition of analytic multiplier
ideals, finite normalization of an analytic blowup, and
generic full rank of a finite map in characteristic zero.
The determinant estimate (4) and its application to (2)--(3)
are proved in this paper. The specialization of the
Jacobian--multiplier inclusion is also written out.

The source for that inclusion is L.~Ein, R.~Lazarsfeld,
K.~E.~Smith and D.~Varolin, \emph{Jumping coefficients of
multiplier ideals}, Duke Math.~J.~123(3) (2004), 469--506,
Definition4.1, Theorem4.2(i) and Lemma4.3.
The pertinent complete proof was read in the
\href{https://www.math.stonybrook.edu/~dror/paper-jump.pdf}
{author-hosted 29-page text}, Section4.
The discrepancy application is compared with
\href{https://arxiv.org/pdf/2605.30063v2}{Trusiani2605.30063v2},
Proposition3.1 and AppendixA.1, and
\href{https://arxiv.org/pdf/2609.31089v1}{Trusiani2609.31089v1},
Proposition4.1.
The source versions and proof applications are distinguished;
no third-party PDF or lengthy excerpt is reproduced.

\end{document}
